\documentclass[11pt]{article}

\usepackage[T1]{fontenc}
\usepackage{lmodern}
\usepackage{amsmath}
\usepackage{amssymb}
\usepackage[margin=1in]{geometry}
\usepackage{booktabs}
\usepackage{hyperref}

\newcommand{\key}[1]{\texttt{#1}}

\title{Setting the load-balancing parameters of the Monte Carlo module}
\author{Athena++ Monte Carlo radiation transport}
\date{}

\begin{document}
\sloppy
\maketitle

\section{What is balanced}

Every block records the wall time its transport sweeps take, and the module adds that
time to the mesh's per-block cost. The mesh's balancer then redistributes blocks across
ranks so that the predicted cost of the busiest rank is as small as possible, and the
module moves the photons resident on a block, its moment arrays and its random-number
state along with it. Costs are measured, not modelled, so the first transport of a run
starts with whatever layout the mesh built, and every later decision uses what the
previous transport actually cost.

Two things follow. The layout can only improve once a transport has been timed, so a
run of a single output interval gains nothing unless a cost file from an earlier run or
the photon-count guess described below is supplied. And costs are only comparable when
the work per photon is stationary, which holds for a static run over its output
intervals and for the repeated transports of a dynamic run.

\section{The two cadences}

Redistribution can happen at two points, controlled by keys in two blocks of the deck.

\paragraph{Between output intervals.}
The mesh's own machinery, with keys in \key{<loadbalancing>}. At the start of each
\key{RunMonteCarlo} call the mesh compares the most expensive rank against the mean, and
if \key{interval} transports have passed since the last check and the busiest rank costs
more than $(1+\key{tolerance})$ times the average, it repartitions. In a static run each
output interval is one transport, so \key{interval = 1} considers a redistribution at
every interval with the costs of the interval just finished. The first interval of a
static run is unbalanced unless \key{cost\_file} or \key{lb\_initial} provides costs.

\paragraph{Mid-transport.}
Enabled by \key{<montecarlo>/lb\_check\_interval} $> 0$. Which transport loop is running
decides what a check means:
\begin{itemize}
  \item In the synchronous loop (one rank, or the asynchronous termination switched
    off), a check happens every \key{lb\_check\_interval} completed rounds, and at
    least \key{lb\_min\_window} rounds after the last redistribution.
  \item In the asynchronous loop, the production MPI case, a check happens each time
    \key{lb\_check\_fraction} of the interval's photons have finished since the last
    check. \key{lb\_check\_interval} only has to be nonzero to enable it; its value is
    not used as a count.
\end{itemize}
In both loops a check gathers the block costs, and every rank evaluates the same
partition. A redistribution is taken only if the predicted busiest rank improves by at
least \key{lb\_min\_gain}, and at most \key{lb\_max\_per\_transport} redistributions are
made in one transport. In the asynchronous loop the ranks first drain: they stop
transporting, finish delivering photons in flight, and make the blocking balance call
only when every rank reports the same sent and received totals twice in a row.

\paragraph{With \key{nout} $> 1$.}
All the mid-transport counters are local to one transport, and the fraction is taken of
that interval's \key{nphot}, not of \key{nout} $\times$ \key{nphot}. A run with
\key{nout = 3} and \key{lb\_check\_fraction = 0.1} checks after every ten per cent of
each interval's photons, up to \key{lb\_max\_per\_transport} times in each interval, and
additionally considers a between-interval balance at the top of the second and third
intervals. What persists across intervals is the cost information and the layout.

\section{The keys}

\begin{center}
\small
\begin{tabular}{@{}lp{3.4cm}p{7.4cm}@{}}
\toprule
Key & Default & Meaning \\
\midrule
\multicolumn{3}{@{}l}{\key{<loadbalancing>}} \\
\key{balancer} & \key{default} & \key{automatic} balances on measured cost;
  \key{manual} uses the cost list the module sets. \key{default} turns balancing off. \\
\key{tolerance} & 0.5 & A between-interval redistribution needs the busiest rank
  above $(1+\key{tolerance})$ times the mean. \\
\key{interval} & 10 & Transports between between-interval checks. Set 1 for a static
  run with several outputs. \\
\key{cost\_file} & none & Read block costs from this file before the first transport and
  write them after each transport, so a rerun starts balanced. \\
\midrule
\multicolumn{3}{@{}l}{\key{<montecarlo>}} \\
\key{lb\_initial} & \key{none} & \key{photons}: with no costs on hand, balance the
  first transport on each block's share of the photons to emit. Useful with
  \key{weights = equal} or \key{biased}, where that share follows the emission; with
  \key{weights = emission} every block gets the same count and it does nothing. \\
\key{lb\_check\_interval} & 0 & Enables mid-transport checks; in the synchronous loop,
  the number of rounds between them. \\
\key{lb\_check\_fraction} & 0.1 & Asynchronous loop: fraction of the interval's photons
  finished between checks. \\
\key{lb\_max\_per\_transport} & 4 & Cap on redistributions within one transport. \\
\key{lb\_min\_window} & \key{lb\_check\_interval} & Synchronous loop: rounds that must
  pass after a redistribution before the next check. \\
\key{lb\_min\_gain} & 0.05 & Required fractional improvement of the predicted busiest
  rank for a redistribution to be taken, between intervals and mid-transport. \\
\key{lb\_partition} & \key{optimal} & Contiguous partition of the block list:
  \key{optimal} minimises the busiest rank exactly; \key{greedy} is the mesh's
  cheaper heuristic. \\
\key{lb\_cost\_decay} & 1.0 & After a check that keeps the layout, multiply the
  accumulated costs by this. Below 1 it forgets old cost so a change in where the work
  is shows up sooner. \\
\key{lb\_report} & \key{false} & Print the per-rank and per-block cost report after each
  transport and the wall cost of each redistribution. \\
\key{lb\_test\_repack} \newline \key{lb\_test\_costs} & \key{none} \newline \key{measured} & Test
  knobs that force moves or substitute costs. Leave at their defaults in production. \\
\bottomrule
\end{tabular}
\normalsize
\end{center}

\section{Recommended settings}

\paragraph{Static run, several output intervals.}
\begin{verbatim}
<loadbalancing>
balancer  = automatic
tolerance = 0.2
interval  = 1
cost_file = costs.dat

<montecarlo>
lb_initial        = photons
lb_check_interval = 1
lb_check_fraction = 0.1
lb_report         = true
\end{verbatim}
The first interval runs on the photon-count guess, or on the cost file if one exists from
an earlier run of the same snapshot; the mid-transport checks catch the imbalance
within that interval; every later interval starts from the measured costs. Keep
\key{lb\_report} on until the report shows the busiest rank within a few tens of per cent
of the fair share, then turn it off. \key{tolerance} at 0.5 is the hydro default and is
loose for a transport whose cost is dominated by a few blocks; 0.2 makes the
between-interval check act on the imbalance the report typically shows.

\paragraph{Single output interval.}
Between-interval balancing never fires. Use \key{cost\_file} from a previous run of the
same snapshot if there is one, otherwise \key{lb\_initial = photons} together with the
mid-transport checks. Raise \key{lb\_max\_per\_transport} only if the report shows the
cap being reached while the busiest rank is still far above the fair share.

\paragraph{Dynamic run.}
Each cycle is a transport, so \key{interval} counts cycles, and the costs accumulate
from cycle to cycle unless \key{lb\_cost\_decay} is below 1. The mid-transport
machinery is not used in dynamic runs; the between-interval keys are what matter.

\section{What to expect and what to watch}

\begin{itemize}
  \item Redistribution moves a block's random-number state with it, so a balanced run
    is not bitwise equal to an unbalanced one; spectra agree within the Poisson noise,
    and photon counts are conserved exactly. The report prints the sent and received
    totals of each move.
  \item A move costs a collective and the transfer of every resident photon on the
    blocks that move. The prediction guard exists so that a move is not made for a
    gain smaller than its cost; if the report shows moves with small gains,
    raise \key{lb\_min\_gain}.
  \item The partition is contiguous in the block list, so a single very expensive block
    cannot be split; the best achievable busiest rank is at least that block's cost.
    The per-block part of the report shows whether the imbalance is one block or many.
  \item The \key{mc\_readhdf} and \key{mc\_readhdf\_gr} generators refuse balancing
    unless \key{emission = freefree}, because their per-block opacity tables are not
    rebuilt when a block moves. The free-free path rebuilds its number densities on
    arrival and is safe.
  \item With \key{weights = biased} the importance and escape tables travel with the
    block, so balancing is safe in that scheme too.
\end{itemize}

\end{document}
